\section{Counterfactual Evaluation Framework}
\subsection{Reference sensitivity before method evaluation}
Exact/exact comparisons define an empirical null envelope. Original/exact discrepancy is
then compared with that envelope before any approximate-method decision. This ordering is
essential: when deleting $D_f$ does not create a detectable shift under metric $j$, a low
approximate/exact distance cannot identify whether $U$ removed influence or merely
preserved behavior that was already compatible with retraining.

The P001 retrospective analysis summarizes this relation with a separability ratio formed
from the original/exact signal and the exact/exact reference scale. Power thresholds were
simulated with seed-level summaries. Because this analysis was designed after P001, it is
reported as exploratory and cannot retroactively establish a confirmatory threshold.

\subsection{Fidelity dimensions}
The full framework distinguishes the following coordinates:
\begin{description}
\item[Predictive fidelity.] Agreement, disagreement, and class-conditioned errors.
\item[Probabilistic fidelity.] Distances between complete predictive distributions,
including JS divergence and total variation.
\item[Calibration fidelity.] NLL, Brier score, ECE, and reliability behavior on forgotten,
retained, and independent test sets.
\item[Privacy fidelity.] Difference between attacks against $M_U$ and attacks against the
exact population, rather than comparison with chance alone.
\item[Representation fidelity.] CKA or other symmetry-aware comparisons on a fixed probe
set for compatible architectures.
\item[Local and robustness fidelity.] Behavioral change near deleted observations and under
controlled corruptions, without mislabeling corruption tests as adversarial robustness.
\item[Utility and cost.] Retained/test performance, runtime, device, memory, and cost
relative to exact retraining.
\end{description}

\subsection{Decision states}
Composite rules and equivalence margins must be frozen before confirmatory experiments.
The current vocabulary distinguishes assay-insensitive (U0), functional failure (U1),
functional-only equivalence (U2), strong empirical equivalence (U3), and inconclusive
(U4). U3 would remain an empirical, assay-bounded statement rather than certified removal
or a legal claim. P001 is assigned U0 retrospectively for its two evaluated output probes;
no method receives a confirmatory state.
